\section{Boolean Variable Representation} \label{sec:boolean-variable-representation}
The encoding uses several families of propositional variables.
The most important ones are shown in Table~\ref{tab:variables}.

\begin{table}[H]
    \centering
    \begin{tabular}{ll}
        \toprule
        \textbf{Variable}         & \textbf{Meaning}                              \\
        \midrule
        $\texttt{atFloor}(p,f,t)$ & Person $p$ is on floor $f$ at time $t$        \\
        $\texttt{atFloor}(i,f,t)$ & Item $i$ is on floor $f$ at time $t$          \\
        $\texttt{ascends}(p,t)$   & Person $p$ ascends one floor during step $t$  \\
        $\texttt{descends}(p,t)$  & Person $p$ descends one floor during step $t$ \\
        $\texttt{carries}(p,i,t)$ & Person $p$ carries item $i$ during step $t$   \\
        \bottomrule
    \end{tabular}
    \caption{Boolean variable families used in the SAT encoding.}
    \label{tab:variables}
\end{table}

\noindent In the implementation, each symbolic variable is assigned a unique positive integer using a PySAT \texttt{IDPool}.
This is necessary because SAT solvers operate on integer variable identifiers rather than symbolic names.
A separate mapping is produced so that the integer identifiers can be translated back into readable names such as:
\[
    \texttt{atFloor(Alice, 2, 5)}
\]
or
\[
    \texttt{carries(Bob, Sofa, 7)}.
\]

\noindent For a problem with $n$ floors, $m$ movers, $k$ items, and time horizon $T$, the variable families have the following approximate sizes:

\begin{itemize}
    \item Person position variables: $m \cdot n \cdot (T+1)$.
    \item Item position variables: $k \cdot n \cdot (T+1)$.
    \item Ascend variables: $m \cdot T$.
    \item Descend variables: $m \cdot T$.
    \item Carry variables: $m \cdot k \cdot T$.
\end{itemize}

Therefore, Task 2 uses approximately
\[
    m n (T+1) + k n (T+1) + 2mT + mkT
\]
Boolean variables.